\section{Запрос и ответ службы штампов времени}\label{FMT.TSP}

\subsection{Формат запроса}\label{FMT.TSP.Req}

Формат запроса на получение штампа времени задается типом~\code{TimeStampReq},
который определен в СТБ~34.101.82.

В~\code{TimeStampReq} должны быть опущены опциональные
компоненты~\code{reqPolicy}, \code{nonce} и~\code{extensions}.
%
Если в ответе не требуется получить сертификат СШВ, то должен быть также опущен
компонент~\code{certReq}, который по умолчанию принимает значение~\code{FALSE}.
%
Если сертификат нужен, то компонент~\code{certReq} должен включаться со
значением~\code{TRUE}.

В компоненте~\code{messageImprint} контейнера~\code{TimeStampReq} указывается
хэш-значение данных. Хэш-значение сопровождается описанием алгоритма
хэширования. СШВ должна проверять соответствие алгоритма и длины хэш-значения.
СШВ должна поддерживать по крайней мере алгоритмы \code{belt-hash},
\code{bash384} и~\code{bash512} (см.~\ref{CRYPTO.Hash}).

\subsection{Формат ответа}\label{FMT.TSP.Resp}
 
Формат ответа СШВ задается типом~\code{TimeStampResp}, который определен 
в СТБ~34.101.82. 

Статус обработки запроса указывается в компоненте~\code{status}.
%
В случае успеха вложенные в~\code{status} компоненты~\code{statusString}
и~\code{failInfo} должны быть опущены, а вложенный компонент~\code{status}
должен принимать значение~\code{granted}.
%
Если ошибка произошла по причине загруженности сервера, то~\code{statusString}
и~\code{failInfo} также должны быть опущены, а~\code{status} должен принимать
значение~\code{waiting}.
%
При других ошибках в~\code{status} должно быть выбрано
значение~\code{rejection}, компонент~\code{statusString} должен быть опущен, а
в~\code{failInfo} должен быть указан один из кодов ошибки, определенных в СТБ
34.101.82.

Штамп времени указывается в компоненте~\code{timeStampToken}. Штамп в свою
очередь содержит контейнер типа~\code{SignedData}.
%
Контейнер~\code{SignedData} должен формироваться по правилам, заданным
в~\ref{FMT.SignedData}. Дополнительное правило  касается
компонента~\code{certificates} с сертификатом СШВ: этот компонент должен
отсутствовать, если в запросе на получение штампа времени опущен
флаг~\code{certReq}.

Контейнер~\code{SignedData} должен содержать ссылку на сертификат СШВ в виде
подписанного атрибута~\code{SigningCertificate} или~\code{SigningCertificateV2}.
Эти атрибуты определены в СТБ~34.101.80.
%
В~\code{SigningCertificate} должен быть опущен компонент~\code{policies}, должен
быть вложен только один контейнер~\code{ESSCertID} (ссылка на сертификат СШВ), и
в этом контейнере должен быть опущен компонент~\code{issuerSerial}.
%
Аналогично: в~\code{SigningCertificateV2} должен быть опущен
компонент~\code{policies}, в единственном вложенном
контейнере~\code{ESSCertIDv2}~--- \code{issuerSerial}.

В~\code{SignedData} непосредственно подписывается контейнер типа~\code{TSTInfo}.
%
В подписываемом контейнере должны быть опущены все опциональные компоненты,
кроме, возможно, \code{accuracy}.
%
В компоненте~\code{version} должно быть установлено значение~$1$, а в
компоненте~\code{policy}~--- значение~\code{bpki-role-tsa}, определенное в
приложении~\ref{ASN1}.
%
В компонент~\code{messageImprint} должно быть перенесено значение одноименного
компонента запроса.
